% --- Archon physics formalization source begin ---
% archon:physics
% archon:covers PhyXMiniProblems/problem_phyx_mini_0938.lean
% archon:source-report reports/phyx_mini/problem_phyx_mini_0938.source.json

\chapter{Physics problem phyx\_mini\_0938}
\label{ch:PhyXMiniProblems_problem_phyx_mini_0938}

\paragraph{Problem source.}
\#\# Physical scenario

Suppose the long solenoid in figure has 500 turns per meter and cross-sectional area \textbackslash{}( 4.0 \textbackslash{}\textbackslash{} \textbackslash{}mathrm\{cm\}\textasciicircum{}2 \textbackslash{}). The current in its windings is increasing at \textbackslash{}( 100 \textbackslash{}\textbackslash{} \textbackslash{}mathrm\{A/s\} \textbackslash{}).

\#\# Question

Find the magnitude of the induced electric field within the loop if its radius is \textbackslash{}( 2.0 \textbackslash{}\textbackslash{} \textbackslash{}mathrm\{cm\} \textbackslash{}).

\#\# Answer choices

- A: A: 2.2 \textbackslash{}times 10\textasciicircum{}\{-4\}V/m

- B: B: 2 \textbackslash{}times 10\textasciicircum{}\{-4\}V/m

- C: C: 1 \textbackslash{}times 10\textasciicircum{}\{-4\}V/m

- D: D: 2.6 \textbackslash{}times 10\textasciicircum{}\{-4\}V/m

\#\# Figure caption (auxiliary; use the image as primary evidence)

The image depicts a solenoid with a wire loop surrounding it. Key features include:

1. **Solenoid**: A cylindrical structure wound with a black coil, labeled as "Solenoid." The solenoid carries a current \textbackslash{}( I \textbackslash{}), changing at a rate \textbackslash{}( \textbackslash{}frac\{dI\}\{dt\} \textbackslash{}).

2. **Blue Cylinder**: Represents the region with a magnetic field, denoted as \textbackslash{}( \textbackslash{}vec\{B\} \textbackslash{}). The magnetic field direction is indicated by a blue arrow along the solenoid.

3. **Wire Loop**: Encircles the solenoid and is connected to a galvanometer, labeled with "G." This loop carries an induced current \textbackslash{}( I' \textbackslash{}).

4. **Galvanometer**: Measures the induced current in the wire loop. 

5. **Labeling**: The directions of both \textbackslash{}( I \textbackslash{}) and \textbackslash{}( \textbackslash{}frac\{dI\}\{dt\} \textbackslash{}) are marked with purple arrows, indicating the current direction and its rate of change.

6. **Text**: Annotations are present describing the solenoid, wire loop, galvanometer, and the magnetic field region within the blue cylinder.

\#\# Dataset metadata

Electromagnetic Induction; Physical Model Grounding Reasoning, Multi-Formula Reasoning

\paragraph{Recorded answer/context.}
B: B: 2 \textbackslash{}times 10\textasciicircum{}\{-4\}V/m

\paragraph{Figure/image path.}
/root/proposal\_for\_physic/science-mango/phyx\_mini\_run/phyx\_data/test\_image/938.png

\paragraph{Formalization target.}
create a compiling Lean file with sorry bodies at `PhyXMiniProblems/problem\_phyx\_mini\_0938.lean`. The Lean declarations must preserve the physical quantities, dimensions or dimensional roles, figure labels, governing-law hypotheses, and final relation expressed by this problem.
Use Mathlib/Physlib names found through LeanExplore where available. If a physics API is missing, introduce faithful local abstractions rather than scalar placeholder aliases.

\begin{theorem}[Physics formalization target]
\label{thm:physics:phyx_mini_0938:target}
\lean{PhyXMiniProblems.ProblemPhyXMini0938.problem_phyx_mini_0938}
\uses{def:physics:phyx-mini-0938:phyxminiproblems-problemphyxmini0938-electricfieldmagnitudeinvoltspermeter, def:physics:phyx-mini-0938:phyxminiproblems-problemphyxmini0938-solenoidinductionsetup, def:physics:phyx-mini-0938:phyxminiproblems-problemphyxmini0938-matcheslongsolenoidloopscenario, def:physics:phyx-mini-0938:phyxminiproblems-problemphyxmini0938-matchessuppliedsolenoidloopfigure, def:physics:phyx-mini-0938:phyxminiproblems-problemphyxmini0938-hasphysicalsolenoidloopparameters, def:physics:phyx-mini-0938:phyxminiproblems-problemphyxmini0938-hasstatedsolenoidloopdata, def:physics:phyx-mini-0938:phyxminiproblems-problemphyxmini0938-hastextbookvacuumpermeability, def:physics:phyx-mini-0938:phyxminiproblems-problemphyxmini0938-satisfiescircularloopgeometry, def:physics:phyx-mini-0938:phyxminiproblems-problemphyxmini0938-satisfieslongsolenoidfieldlaw, def:physics:phyx-mini-0938:phyxminiproblems-problemphyxmini0938-satisfiesuniformsolenoidfluxlaw, def:physics:phyx-mini-0938:phyxminiproblems-problemphyxmini0938-satisfiesfaradayinductionlaw, def:physics:phyx-mini-0938:phyxminiproblems-problemphyxmini0938-satisfiescircularinducedelectricfieldlaw, def:physics:phyx-mini-0938:phyxminiproblems-problemphyxmini0938-answerchoice, def:physics:phyx-mini-0938:phyxminiproblems-problemphyxmini0938-recordeddatasetanswer}
The assigned autoformalize agent should translate this physics problem into Lean declarations in the covered file, with theorem and lemma proof bodies written as `by sorry`.
\end{theorem}
\begin{proof}
This is an autoformalization task, not a proof task. Produce statements that can later be proved by the physics prover without weakening the physical model.
\end{proof}

% --- Archon declaration topology begin ---
\section{Declaration topology}
\begin{definition}[electric Current Dimension]
  \label{def:physics:phyx-mini-0938:phyxminiproblems-problemphyxmini0938-electriccurrentdimension}
  \lean{PhyXMiniProblems.ProblemPhyXMini0938.electricCurrentDimension}
  Electric current has dimension C T⁻¹
\end{definition}
\begin{proof}
  This is the declared model definition of electric Current Dimension.
\end{proof}
\begin{definition}[electric Current Rate Dimension]
  \label{def:physics:phyx-mini-0938:phyxminiproblems-problemphyxmini0938-electriccurrentratedimension}
  \lean{PhyXMiniProblems.ProblemPhyXMini0938.electricCurrentRateDimension}
  \uses{def:physics:phyx-mini-0938:phyxminiproblems-problemphyxmini0938-electriccurrentdimension}
  A current-change rate has dimension C T⁻²
\end{definition}
\begin{proof}
  This is the declared model definition of electric Current Rate Dimension.
\end{proof}
\begin{definition}[turn Density Dimension]
  \label{def:physics:phyx-mini-0938:phyxminiproblems-problemphyxmini0938-turndensitydimension}
  \lean{PhyXMiniProblems.ProblemPhyXMini0938.turnDensityDimension}
  Turns per unit length have inverse-length dimension L⁻¹
\end{definition}
\begin{proof}
  This is the declared model definition of turn Density Dimension.
\end{proof}
\begin{definition}[area Dimension]
  \label{def:physics:phyx-mini-0938:phyxminiproblems-problemphyxmini0938-areadimension}
  \lean{PhyXMiniProblems.ProblemPhyXMini0938.areaDimension}
  Area has dimension L²
\end{definition}
\begin{proof}
  This is the declared model definition of area Dimension.
\end{proof}
\begin{definition}[magnetic Flux Density Dimension]
  \label{def:physics:phyx-mini-0938:phyxminiproblems-problemphyxmini0938-magneticfluxdensitydimension}
  \lean{PhyXMiniProblems.ProblemPhyXMini0938.magneticFluxDensityDimension}
  Magnetic flux density has dimension M T⁻¹ C⁻¹
\end{definition}
\begin{proof}
  This is the declared model definition of magnetic Flux Density Dimension.
\end{proof}
\begin{definition}[magnetic Flux Density Rate Dimension]
  \label{def:physics:phyx-mini-0938:phyxminiproblems-problemphyxmini0938-magneticfluxdensityratedimension}
  \lean{PhyXMiniProblems.ProblemPhyXMini0938.magneticFluxDensityRateDimension}
  \uses{def:physics:phyx-mini-0938:phyxminiproblems-problemphyxmini0938-magneticfluxdensitydimension}
  A magnetic-flux-density change rate has dimension M T⁻² C⁻¹
\end{definition}
\begin{proof}
  This is the declared model definition of magnetic Flux Density Rate Dimension.
\end{proof}
\begin{definition}[magnetic Flux Dimension]
  \label{def:physics:phyx-mini-0938:phyxminiproblems-problemphyxmini0938-magneticfluxdimension}
  \lean{PhyXMiniProblems.ProblemPhyXMini0938.magneticFluxDimension}
  \uses{def:physics:phyx-mini-0938:phyxminiproblems-problemphyxmini0938-areadimension, def:physics:phyx-mini-0938:phyxminiproblems-problemphyxmini0938-magneticfluxdensitydimension}
  Magnetic flux has dimension M L² T⁻¹ C⁻¹
\end{definition}
\begin{proof}
  This is the declared model definition of magnetic Flux Dimension.
\end{proof}
\begin{definition}[electromotive Force Dimension]
  \label{def:physics:phyx-mini-0938:phyxminiproblems-problemphyxmini0938-electromotiveforcedimension}
  \lean{PhyXMiniProblems.ProblemPhyXMini0938.electromotiveForceDimension}
  \uses{def:physics:phyx-mini-0938:phyxminiproblems-problemphyxmini0938-magneticfluxdimension}
  Magnetic-flux rate and electromotive force have dimension M L² T⁻² C⁻¹
\end{definition}
\begin{proof}
  This is the declared model definition of electromotive Force Dimension.
\end{proof}
\begin{definition}[electric Field Dimension]
  \label{def:physics:phyx-mini-0938:phyxminiproblems-problemphyxmini0938-electricfielddimension}
  \lean{PhyXMiniProblems.ProblemPhyXMini0938.electricFieldDimension}
  \uses{def:physics:phyx-mini-0938:phyxminiproblems-problemphyxmini0938-electromotiveforcedimension}
  Electric field has dimension M L T⁻² C⁻¹, equivalently volts per metre
\end{definition}
\begin{proof}
  This is the declared model definition of electric Field Dimension.
\end{proof}
\begin{definition}[Length Quantity]
  \label{def:physics:phyx-mini-0938:phyxminiproblems-problemphyxmini0938-lengthquantity}
  \lean{PhyXMiniProblems.ProblemPhyXMini0938.LengthQuantity}
  A nonnegative, unit-independent physical length
\end{definition}
\begin{proof}
  This is the declared model definition of Length Quantity.
\end{proof}
\begin{definition}[Area Quantity]
  \label{def:physics:phyx-mini-0938:phyxminiproblems-problemphyxmini0938-areaquantity}
  \lean{PhyXMiniProblems.ProblemPhyXMini0938.AreaQuantity}
  \uses{def:physics:phyx-mini-0938:phyxminiproblems-problemphyxmini0938-areadimension}
  A nonnegative, unit-independent area
\end{definition}
\begin{proof}
  This is the declared model definition of Area Quantity.
\end{proof}
\begin{definition}[Current Magnitude Quantity]
  \label{def:physics:phyx-mini-0938:phyxminiproblems-problemphyxmini0938-currentmagnitudequantity}
  \lean{PhyXMiniProblems.ProblemPhyXMini0938.CurrentMagnitudeQuantity}
  \uses{def:physics:phyx-mini-0938:phyxminiproblems-problemphyxmini0938-electriccurrentdimension}
  A nonnegative winding-current magnitude
\end{definition}
\begin{proof}
  This is the declared model definition of Current Magnitude Quantity.
\end{proof}
\begin{definition}[Current Increase Rate Quantity]
  \label{def:physics:phyx-mini-0938:phyxminiproblems-problemphyxmini0938-currentincreaseratequantity}
  \lean{PhyXMiniProblems.ProblemPhyXMini0938.CurrentIncreaseRateQuantity}
  \uses{def:physics:phyx-mini-0938:phyxminiproblems-problemphyxmini0938-electriccurrentratedimension}
  A nonnegative winding-current increase rate
\end{definition}
\begin{proof}
  This is the declared model definition of Current Increase Rate Quantity.
\end{proof}
\begin{definition}[Turn Density Quantity]
  \label{def:physics:phyx-mini-0938:phyxminiproblems-problemphyxmini0938-turndensityquantity}
  \lean{PhyXMiniProblems.ProblemPhyXMini0938.TurnDensityQuantity}
  \uses{def:physics:phyx-mini-0938:phyxminiproblems-problemphyxmini0938-turndensitydimension}
  A nonnegative number of solenoid turns per unit length
\end{definition}
\begin{proof}
  This is the declared model definition of Turn Density Quantity.
\end{proof}
\begin{definition}[Magnetic Flux Density Quantity]
  \label{def:physics:phyx-mini-0938:phyxminiproblems-problemphyxmini0938-magneticfluxdensityquantity}
  \lean{PhyXMiniProblems.ProblemPhyXMini0938.MagneticFluxDensityQuantity}
  \uses{def:physics:phyx-mini-0938:phyxminiproblems-problemphyxmini0938-magneticfluxdensitydimension}
  A nonnegative magnetic-flux-density magnitude
\end{definition}
\begin{proof}
  This is the declared model definition of Magnetic Flux Density Quantity.
\end{proof}
\begin{definition}[Magnetic Flux Density Increase Rate Quantity]
  \label{def:physics:phyx-mini-0938:phyxminiproblems-problemphyxmini0938-magneticfluxdensityincreaseratequantity}
  \lean{PhyXMiniProblems.ProblemPhyXMini0938.MagneticFluxDensityIncreaseRateQuantity}
  \uses{def:physics:phyx-mini-0938:phyxminiproblems-problemphyxmini0938-magneticfluxdensityratedimension}
  A nonnegative magnetic-flux-density increase rate
\end{definition}
\begin{proof}
  This is the declared model definition of Magnetic Flux Density Increase Rate Quantity.
\end{proof}
\begin{definition}[Signed Magnetic Flux Rate Quantity]
  \label{def:physics:phyx-mini-0938:phyxminiproblems-problemphyxmini0938-signedmagneticfluxratequantity}
  \lean{PhyXMiniProblems.ProblemPhyXMini0938.SignedMagneticFluxRateQuantity}
  \uses{def:physics:phyx-mini-0938:phyxminiproblems-problemphyxmini0938-electromotiveforcedimension}
  Signed magnetic-flux change rate relative to the selected area normal
\end{definition}
\begin{proof}
  This is the declared model definition of Signed Magnetic Flux Rate Quantity.
\end{proof}
\begin{definition}[Signed Electric Circulation Quantity]
  \label{def:physics:phyx-mini-0938:phyxminiproblems-problemphyxmini0938-signedelectriccirculationquantity}
  \lean{PhyXMiniProblems.ProblemPhyXMini0938.SignedElectricCirculationQuantity}
  \uses{def:physics:phyx-mini-0938:phyxminiproblems-problemphyxmini0938-electromotiveforcedimension}
  Signed electric-field circulation relative to the selected loop orientation
\end{definition}
\begin{proof}
  This is the declared model definition of Signed Electric Circulation Quantity.
\end{proof}
\begin{definition}[Signed Electric Field Quantity]
  \label{def:physics:phyx-mini-0938:phyxminiproblems-problemphyxmini0938-signedelectricfieldquantity}
  \lean{PhyXMiniProblems.ProblemPhyXMini0938.SignedElectricFieldQuantity}
  \uses{def:physics:phyx-mini-0938:phyxminiproblems-problemphyxmini0938-electricfielddimension}
  Signed tangential component of electric field on the circular loop
\end{definition}
\begin{proof}
  This is the declared model definition of Signed Electric Field Quantity.
\end{proof}
\begin{definition}[Electric Field Magnitude Quantity]
  \label{def:physics:phyx-mini-0938:phyxminiproblems-problemphyxmini0938-electricfieldmagnitudequantity}
  \lean{PhyXMiniProblems.ProblemPhyXMini0938.ElectricFieldMagnitudeQuantity}
  \uses{def:physics:phyx-mini-0938:phyxminiproblems-problemphyxmini0938-electricfielddimension}
  Nonnegative magnitude of electric field on the circular loop
\end{definition}
\begin{proof}
  This is the declared model definition of Electric Field Magnitude Quantity.
\end{proof}
\begin{definition}[length In Meters]
  \label{def:physics:phyx-mini-0938:phyxminiproblems-problemphyxmini0938-lengthinmeters}
  \lean{PhyXMiniProblems.ProblemPhyXMini0938.lengthInMeters}
  \uses{def:physics:phyx-mini-0938:phyxminiproblems-problemphyxmini0938-lengthquantity}
  Read a length in coherent-SI metres
\end{definition}
\begin{proof}
  This is the declared model definition of length In Meters.
\end{proof}
\begin{definition}[area In Square Meters]
  \label{def:physics:phyx-mini-0938:phyxminiproblems-problemphyxmini0938-areainsquaremeters}
  \lean{PhyXMiniProblems.ProblemPhyXMini0938.areaInSquareMeters}
  \uses{def:physics:phyx-mini-0938:phyxminiproblems-problemphyxmini0938-areaquantity}
  Read an area in coherent-SI square metres
\end{definition}
\begin{proof}
  This is the declared model definition of area In Square Meters.
\end{proof}
\begin{definition}[current In Amperes]
  \label{def:physics:phyx-mini-0938:phyxminiproblems-problemphyxmini0938-currentinamperes}
  \lean{PhyXMiniProblems.ProblemPhyXMini0938.currentInAmperes}
  \uses{def:physics:phyx-mini-0938:phyxminiproblems-problemphyxmini0938-currentmagnitudequantity}
  Read a winding current in coherent-SI amperes
\end{definition}
\begin{proof}
  This is the declared model definition of current In Amperes.
\end{proof}
\begin{definition}[current Increase Rate In Amperes Per Second]
  \label{def:physics:phyx-mini-0938:phyxminiproblems-problemphyxmini0938-currentincreaserateinamperespersecond}
  \lean{PhyXMiniProblems.ProblemPhyXMini0938.currentIncreaseRateInAmperesPerSecond}
  \uses{def:physics:phyx-mini-0938:phyxminiproblems-problemphyxmini0938-currentincreaseratequantity}
  Read a winding-current increase rate in coherent-SI amperes per second
\end{definition}
\begin{proof}
  This is the declared model definition of current Increase Rate In Amperes Per Second.
\end{proof}
\begin{definition}[turn Density In Turns Per Meter]
  \label{def:physics:phyx-mini-0938:phyxminiproblems-problemphyxmini0938-turndensityinturnspermeter}
  \lean{PhyXMiniProblems.ProblemPhyXMini0938.turnDensityInTurnsPerMeter}
  \uses{def:physics:phyx-mini-0938:phyxminiproblems-problemphyxmini0938-turndensityquantity}
  Read solenoid turn density in coherent-SI turns per metre
\end{definition}
\begin{proof}
  This is the declared model definition of turn Density In Turns Per Meter.
\end{proof}
\begin{definition}[magnetic Flux Density In Teslas]
  \label{def:physics:phyx-mini-0938:phyxminiproblems-problemphyxmini0938-magneticfluxdensityinteslas}
  \lean{PhyXMiniProblems.ProblemPhyXMini0938.magneticFluxDensityInTeslas}
  \uses{def:physics:phyx-mini-0938:phyxminiproblems-problemphyxmini0938-magneticfluxdensityquantity}
  Read magnetic-flux density in coherent-SI teslas
\end{definition}
\begin{proof}
  This is the declared model definition of magnetic Flux Density In Teslas.
\end{proof}
\begin{definition}[magnetic Flux Density Increase Rate In Teslas Per Second]
  \label{def:physics:phyx-mini-0938:phyxminiproblems-problemphyxmini0938-magneticfluxdensityincreaserateinteslaspersecond}
  \lean{PhyXMiniProblems.ProblemPhyXMini0938.magneticFluxDensityIncreaseRateInTeslasPerSecond}
  \uses{def:physics:phyx-mini-0938:phyxminiproblems-problemphyxmini0938-magneticfluxdensityincreaseratequantity}
  Read magnetic-flux-density increase rate in teslas per second
\end{definition}
\begin{proof}
  This is the declared model definition of magnetic Flux Density Increase Rate In Teslas Per Second.
\end{proof}
\begin{definition}[magnetic Flux Rate In Webers Per Second]
  \label{def:physics:phyx-mini-0938:phyxminiproblems-problemphyxmini0938-magneticfluxrateinweberspersecond}
  \lean{PhyXMiniProblems.ProblemPhyXMini0938.magneticFluxRateInWebersPerSecond}
  \uses{def:physics:phyx-mini-0938:phyxminiproblems-problemphyxmini0938-signedmagneticfluxratequantity}
  Read signed magnetic-flux rate in coherent-SI webers per second
\end{definition}
\begin{proof}
  This is the declared model definition of magnetic Flux Rate In Webers Per Second.
\end{proof}
\begin{definition}[electric Circulation In Volts]
  \label{def:physics:phyx-mini-0938:phyxminiproblems-problemphyxmini0938-electriccirculationinvolts}
  \lean{PhyXMiniProblems.ProblemPhyXMini0938.electricCirculationInVolts}
  \uses{def:physics:phyx-mini-0938:phyxminiproblems-problemphyxmini0938-signedelectriccirculationquantity}
  Read signed electric-field circulation in coherent-SI volts
\end{definition}
\begin{proof}
  This is the declared model definition of electric Circulation In Volts.
\end{proof}
\begin{definition}[signed Electric Field In Volts Per Meter]
  \label{def:physics:phyx-mini-0938:phyxminiproblems-problemphyxmini0938-signedelectricfieldinvoltspermeter}
  \lean{PhyXMiniProblems.ProblemPhyXMini0938.signedElectricFieldInVoltsPerMeter}
  \uses{def:physics:phyx-mini-0938:phyxminiproblems-problemphyxmini0938-signedelectricfieldquantity}
  Read a signed tangential electric field in coherent-SI volts per metre
\end{definition}
\begin{proof}
  This is the declared model definition of signed Electric Field In Volts Per Meter.
\end{proof}
\begin{definition}[electric Field Magnitude In Volts Per Meter]
  \label{def:physics:phyx-mini-0938:phyxminiproblems-problemphyxmini0938-electricfieldmagnitudeinvoltspermeter}
  \lean{PhyXMiniProblems.ProblemPhyXMini0938.electricFieldMagnitudeInVoltsPerMeter}
  \uses{def:physics:phyx-mini-0938:phyxminiproblems-problemphyxmini0938-electricfieldmagnitudequantity}
  Read electric-field magnitude in coherent-SI volts per metre
\end{definition}
\begin{proof}
  This is the declared model definition of electric Field Magnitude In Volts Per Meter.
\end{proof}
\begin{definition}[Direction Vector]
  \label{def:physics:phyx-mini-0938:phyxminiproblems-problemphyxmini0938-directionvector}
  \lean{PhyXMiniProblems.ProblemPhyXMini0938.DirectionVector}
  Dimensionless spatial direction vectors in physical three-space
\end{definition}
\begin{proof}
  This is the declared model definition of Direction Vector.
\end{proof}
\begin{definition}[solenoid Axis Unit Vector]
  \label{def:physics:phyx-mini-0938:phyxminiproblems-problemphyxmini0938-solenoidaxisunitvector}
  \lean{PhyXMiniProblems.ProblemPhyXMini0938.solenoidAxisUnitVector}
  \uses{def:physics:phyx-mini-0938:phyxminiproblems-problemphyxmini0938-directionvector}
  A selected unit vector along the solenoid axis and blue cylinder
\end{definition}
\begin{proof}
  This is the declared model definition of solenoid Axis Unit Vector.
\end{proof}
\begin{definition}[Figure Object]
  \label{def:physics:phyx-mini-0938:phyxminiproblems-problemphyxmini0938-figureobject}
  \lean{PhyXMiniProblems.ProblemPhyXMini0938.FigureObject}
  Physical objects visible in the primary raster 938.png
\end{definition}
\begin{proof}
  This is the declared model definition of Figure Object.
\end{proof}
\begin{definition}[Figure Label]
  \label{def:physics:phyx-mini-0938:phyxminiproblems-problemphyxmini0938-figurelabel}
  \lean{PhyXMiniProblems.ProblemPhyXMini0938.FigureLabel}
  Symbolic or textual labels visible in panel (a)
\end{definition}
\begin{proof}
  This is the declared model definition of Figure Label.
\end{proof}
\begin{definition}[Diagram Direction]
  \label{def:physics:phyx-mini-0938:phyxminiproblems-problemphyxmini0938-diagramdirection}
  \lean{PhyXMiniProblems.ProblemPhyXMini0938.DiagramDirection}
  Qualitative arrow directions in the oblique drawing plane
\end{definition}
\begin{proof}
  This is the declared model definition of Diagram Direction.
\end{proof}
\begin{definition}[Solenoid Loop Figure]
  \label{def:physics:phyx-mini-0938:phyxminiproblems-problemphyxmini0938-solenoidloopfigure}
  \lean{PhyXMiniProblems.ProblemPhyXMini0938.SolenoidLoopFigure}
  \uses{def:physics:phyx-mini-0938:phyxminiproblems-problemphyxmini0938-figureobject, def:physics:phyx-mini-0938:phyxminiproblems-problemphyxmini0938-figurelabel, def:physics:phyx-mini-0938:phyxminiproblems-problemphyxmini0938-diagramdirection}
  Literal, non-numerical presentation data from 938.png
\end{definition}
\begin{proof}
  This is the declared model definition of Solenoid Loop Figure.
\end{proof}
\begin{definition}[Apparatus Model]
  \label{def:physics:phyx-mini-0938:phyxminiproblems-problemphyxmini0938-apparatusmodel}
  \lean{PhyXMiniProblems.ProblemPhyXMini0938.ApparatusModel}
  Idealization selected by the prose description
\end{definition}
\begin{proof}
  This is the declared model definition of Apparatus Model.
\end{proof}
\begin{definition}[Solenoid Induction Setup]
  \label{def:physics:phyx-mini-0938:phyxminiproblems-problemphyxmini0938-solenoidinductionsetup}
  \lean{PhyXMiniProblems.ProblemPhyXMini0938.SolenoidInductionSetup}
  \uses{def:physics:phyx-mini-0938:phyxminiproblems-problemphyxmini0938-lengthquantity, def:physics:phyx-mini-0938:phyxminiproblems-problemphyxmini0938-areaquantity, def:physics:phyx-mini-0938:phyxminiproblems-problemphyxmini0938-currentmagnitudequantity, def:physics:phyx-mini-0938:phyxminiproblems-problemphyxmini0938-currentincreaseratequantity, def:physics:phyx-mini-0938:phyxminiproblems-problemphyxmini0938-turndensityquantity, def:physics:phyx-mini-0938:phyxminiproblems-problemphyxmini0938-magneticfluxdensityquantity, def:physics:phyx-mini-0938:phyxminiproblems-problemphyxmini0938-magneticfluxdensityincreaseratequantity, def:physics:phyx-mini-0938:phyxminiproblems-problemphyxmini0938-signedmagneticfluxratequantity, def:physics:phyx-mini-0938:phyxminiproblems-problemphyxmini0938-signedelectriccirculationquantity, def:physics:phyx-mini-0938:phyxminiproblems-problemphyxmini0938-signedelectricfieldquantity, def:physics:phyx-mini-0938:phyxminiproblems-problemphyxmini0938-electricfieldmagnitudequantity, def:physics:phyx-mini-0938:phyxminiproblems-problemphyxmini0938-directionvector, def:physics:phyx-mini-0938:phyxminiproblems-problemphyxmini0938-solenoidloopfigure, def:physics:phyx-mini-0938:phyxminiproblems-problemphyxmini0938-apparatusmodel}
  Independent physical observables. In particular, the induced electric-field magnitude is not defined from the solenoid parameters or an answer choice
\end{definition}
\begin{proof}
  This is the declared model definition of Solenoid Induction Setup.
\end{proof}
\begin{definition}[Matches Long Solenoid Loop Scenario]
  \label{def:physics:phyx-mini-0938:phyxminiproblems-problemphyxmini0938-matcheslongsolenoidloopscenario}
  \lean{PhyXMiniProblems.ProblemPhyXMini0938.MatchesLongSolenoidLoopScenario}
  \uses{def:physics:phyx-mini-0938:phyxminiproblems-problemphyxmini0938-solenoidaxisunitvector, def:physics:phyx-mini-0938:phyxminiproblems-problemphyxmini0938-solenoidinductionsetup}
  Prose-level apparatus and geometry, excluding all answer-bearing values
\end{definition}
\begin{proof}
  This is the declared model definition of Matches Long Solenoid Loop Scenario.
\end{proof}
\begin{definition}[Matches Supplied Solenoid Loop Figure]
  \label{def:physics:phyx-mini-0938:phyxminiproblems-problemphyxmini0938-matchessuppliedsolenoidloopfigure}
  \lean{PhyXMiniProblems.ProblemPhyXMini0938.MatchesSuppliedSolenoidLoopFigure}
  \uses{def:physics:phyx-mini-0938:phyxminiproblems-problemphyxmini0938-solenoidinductionsetup}
  Literal object, label, and arrow evidence read from the primary raster
\end{definition}
\begin{proof}
  This is the declared model definition of Matches Supplied Solenoid Loop Figure.
\end{proof}
\begin{definition}[Has Physical Solenoid Loop Parameters]
  \label{def:physics:phyx-mini-0938:phyxminiproblems-problemphyxmini0938-hasphysicalsolenoidloopparameters}
  \lean{PhyXMiniProblems.ProblemPhyXMini0938.HasPhysicalSolenoidLoopParameters}
  \uses{def:physics:phyx-mini-0938:phyxminiproblems-problemphyxmini0938-lengthinmeters, def:physics:phyx-mini-0938:phyxminiproblems-problemphyxmini0938-areainsquaremeters, def:physics:phyx-mini-0938:phyxminiproblems-problemphyxmini0938-currentincreaserateinamperespersecond, def:physics:phyx-mini-0938:phyxminiproblems-problemphyxmini0938-turndensityinturnspermeter, def:physics:phyx-mini-0938:phyxminiproblems-problemphyxmini0938-solenoidinductionsetup}
  Positivity and nondegeneracy of the independent physical parameters
\end{definition}
\begin{proof}
  This is the declared model definition of Has Physical Solenoid Loop Parameters.
\end{proof}
\begin{definition}[Has Stated Solenoid Loop Data]
  \label{def:physics:phyx-mini-0938:phyxminiproblems-problemphyxmini0938-hasstatedsolenoidloopdata}
  \lean{PhyXMiniProblems.ProblemPhyXMini0938.HasStatedSolenoidLoopData}
  \uses{def:physics:phyx-mini-0938:phyxminiproblems-problemphyxmini0938-lengthinmeters, def:physics:phyx-mini-0938:phyxminiproblems-problemphyxmini0938-areainsquaremeters, def:physics:phyx-mini-0938:phyxminiproblems-problemphyxmini0938-currentincreaserateinamperespersecond, def:physics:phyx-mini-0938:phyxminiproblems-problemphyxmini0938-turndensityinturnspermeter, def:physics:phyx-mini-0938:phyxminiproblems-problemphyxmini0938-solenoidinductionsetup}
  The numerical values supplied by the text, converted at the explicit SI readout boundary: 4 cm² = 4 / 10⁴ m² and 2 cm = 2 / 100 m
\end{definition}
\begin{proof}
  This is the declared model definition of Has Stated Solenoid Loop Data.
\end{proof}
\begin{definition}[Has Textbook Vacuum Permeability]
  \label{def:physics:phyx-mini-0938:phyxminiproblems-problemphyxmini0938-hastextbookvacuumpermeability}
  \lean{PhyXMiniProblems.ProblemPhyXMini0938.HasTextbookVacuumPermeability}
  \uses{def:physics:phyx-mini-0938:phyxminiproblems-problemphyxmini0938-solenoidinductionsetup}
  Textbook coherent-SI calibration of vacuum permeability
\end{definition}
\begin{proof}
  This is the declared model definition of Has Textbook Vacuum Permeability.
\end{proof}
\begin{definition}[Satisfies Circular Loop Geometry]
  \label{def:physics:phyx-mini-0938:phyxminiproblems-problemphyxmini0938-satisfiescircularloopgeometry}
  \lean{PhyXMiniProblems.ProblemPhyXMini0938.SatisfiesCircularLoopGeometry}
  \uses{def:physics:phyx-mini-0938:phyxminiproblems-problemphyxmini0938-lengthinmeters, def:physics:phyx-mini-0938:phyxminiproblems-problemphyxmini0938-solenoidinductionsetup}
  Circular-loop geometry. This relation mentions neither electric nor magnetic fields and does not choose an answer value
\end{definition}
\begin{proof}
  This is the declared model definition of Satisfies Circular Loop Geometry.
\end{proof}
\begin{definition}[Satisfies Long Solenoid Field Law]
  \label{def:physics:phyx-mini-0938:phyxminiproblems-problemphyxmini0938-satisfieslongsolenoidfieldlaw}
  \lean{PhyXMiniProblems.ProblemPhyXMini0938.SatisfiesLongSolenoidFieldLaw}
  \uses{def:physics:phyx-mini-0938:phyxminiproblems-problemphyxmini0938-currentincreaserateinamperespersecond, def:physics:phyx-mini-0938:phyxminiproblems-problemphyxmini0938-turndensityinturnspermeter, def:physics:phyx-mini-0938:phyxminiproblems-problemphyxmini0938-magneticfluxdensityinteslas, def:physics:phyx-mini-0938:phyxminiproblems-problemphyxmini0938-magneticfluxdensityincreaserateinteslaspersecond, def:physics:phyx-mini-0938:phyxminiproblems-problemphyxmini0938-solenoidinductionsetup}
  Ideal long-solenoid law. The vector field is uniform inside the blue cylinder and negligible outside; its scalar increase rate is μ₀ n dI/dt. This law does not mention flux, circulation, loop radius, or induced electric field
\end{definition}
\begin{proof}
  This is the declared model definition of Satisfies Long Solenoid Field Law.
\end{proof}
\begin{definition}[Satisfies Uniform Solenoid Flux Law]
  \label{def:physics:phyx-mini-0938:phyxminiproblems-problemphyxmini0938-satisfiesuniformsolenoidfluxlaw}
  \lean{PhyXMiniProblems.ProblemPhyXMini0938.SatisfiesUniformSolenoidFluxLaw}
  \uses{def:physics:phyx-mini-0938:phyxminiproblems-problemphyxmini0938-areainsquaremeters, def:physics:phyx-mini-0938:phyxminiproblems-problemphyxmini0938-magneticfluxdensityincreaserateinteslaspersecond, def:physics:phyx-mini-0938:phyxminiproblems-problemphyxmini0938-magneticfluxrateinweberspersecond, def:physics:phyx-mini-0938:phyxminiproblems-problemphyxmini0938-solenoidinductionsetup}
  Because the wire loop encloses the whole blue solenoid cross-section, the flux rate is the uniform interior field rate times the solenoid area. The law does not mention electric field or loop circumference
\end{definition}
\begin{proof}
  This is the declared model definition of Satisfies Uniform Solenoid Flux Law.
\end{proof}
\begin{definition}[Satisfies Faraday Induction Law]
  \label{def:physics:phyx-mini-0938:phyxminiproblems-problemphyxmini0938-satisfiesfaradayinductionlaw}
  \lean{PhyXMiniProblems.ProblemPhyXMini0938.SatisfiesFaradayInductionLaw}
  \uses{def:physics:phyx-mini-0938:phyxminiproblems-problemphyxmini0938-magneticfluxrateinweberspersecond, def:physics:phyx-mini-0938:phyxminiproblems-problemphyxmini0938-electriccirculationinvolts, def:physics:phyx-mini-0938:phyxminiproblems-problemphyxmini0938-solenoidinductionsetup}
  Faraday's signed law for the chosen positive surface normal and boundary orientation. It does not mention the solenoid parameters or field magnitude
\end{definition}
\begin{proof}
  This is the declared model definition of Satisfies Faraday Induction Law.
\end{proof}
\begin{definition}[Satisfies Circular Induced Electric Field Law]
  \label{def:physics:phyx-mini-0938:phyxminiproblems-problemphyxmini0938-satisfiescircularinducedelectricfieldlaw}
  \lean{PhyXMiniProblems.ProblemPhyXMini0938.SatisfiesCircularInducedElectricFieldLaw}
  \uses{def:physics:phyx-mini-0938:phyxminiproblems-problemphyxmini0938-lengthinmeters, def:physics:phyx-mini-0938:phyxminiproblems-problemphyxmini0938-electriccirculationinvolts, def:physics:phyx-mini-0938:phyxminiproblems-problemphyxmini0938-signedelectricfieldinvoltspermeter, def:physics:phyx-mini-0938:phyxminiproblems-problemphyxmini0938-electricfieldmagnitudeinvoltspermeter, def:physics:phyx-mini-0938:phyxminiproblems-problemphyxmini0938-solenoidinductionsetup}
  Rotational symmetry makes the induced electric field tangential and constant in signed component on the circular path. The circulation is circumference times that component, while the requested magnitude is its absolute value. No numerical answer occurs in this law
\end{definition}
\begin{proof}
  This is the declared model definition of Satisfies Circular Induced Electric Field Law.
\end{proof}
\begin{definition}[Answer Choice]
  \label{def:physics:phyx-mini-0938:phyxminiproblems-problemphyxmini0938-answerchoice}
  \lean{PhyXMiniProblems.ProblemPhyXMini0938.AnswerChoice}
  Answer labels printed with the question
\end{definition}
\begin{proof}
  This is the declared model definition of Answer Choice.
\end{proof}
\begin{definition}[displayed Electric Field In Volts Per Meter]
  \label{def:physics:phyx-mini-0938:phyxminiproblems-problemphyxmini0938-answerchoice-displayedelectricfieldinvoltspermeter}
  \lean{PhyXMiniProblems.ProblemPhyXMini0938.AnswerChoice.displayedElectricFieldInVoltsPerMeter}
  \uses{def:physics:phyx-mini-0938:phyxminiproblems-problemphyxmini0938-answerchoice}
  Displayed electric-field magnitude for each choice, in volts per metre
\end{definition}
\begin{proof}
  This is the declared model definition of displayed Electric Field In Volts Per Meter.
\end{proof}
\begin{definition}[recorded Dataset Answer]
  \label{def:physics:phyx-mini-0938:phyxminiproblems-problemphyxmini0938-recordeddatasetanswer}
  \lean{PhyXMiniProblems.ProblemPhyXMini0938.recordedDatasetAnswer}
  \uses{def:physics:phyx-mini-0938:phyxminiproblems-problemphyxmini0938-answerchoice}
  Answer label recorded by the source dataset
\end{definition}
\begin{proof}
  This is the declared model definition of recorded Dataset Answer.
\end{proof}
% --- Archon declaration topology end ---
% --- Archon physics formalization source end ---
